%% =====================================================================
%%  B3-T-01-03 -- what B3 contributes to "Power as an Amoral Instrument"
%%  -------------------------------------------------------------------
%%  LAYER A: APPEND-ONLY. Written once, then left alone. Material found
%%  only here sits in the `unique` slot and is protected by rule R10.
%% =====================================================================
%% @kind: source
%% @id: B3-T-01-03
%% @book: B3
%% @topic: T-01-03
%% @locator: preface; Laws 18, 20
%% @status: distilled
%% @unique: yes
%% @integrated: yes
%% @updated: 2026-09-19

\dossier{B3}{T-01-03}{\bookshortB3}{preface; Laws 18, 20}

\begin{distilled}
  \item Power is treated descriptively, as an observable set of dynamics, with moral judgement left to the reader's choice of application.
  \item The court is used as the model system: a closed hierarchy where advancement depends on managing perception, dependence and timing.
  \item Isolation is explicitly priced as a strategic error rather than a refuge --- a fortress cuts you off from information and makes you a conspicuous target.
  \item Non-commitment is treated as an available and often superior position: staying unattached preserves optionality and lets others compete for you.
\end{distilled}

\begin{unique}
  \item The court as the canonical model of power dynamics, and the argument that isolation is a cost rather than a protection.
\end{unique}

\begin{terms}
  \item \emph{the court} --- any closed hierarchy in which advancement depends on the perception of superiors rather than on output alone.
  \item \emph{reversal} --- the section of each law stating when the law should be broken; the book's built-in correction against applying a rule mechanically.
\end{terms}
